\chapter{Eigenvectors, Eigenvalues, and Diagonalization}
\label{app:vecteurs-propres-diagonalisation}

Let $V$ be a vector space of finite dimension $N$ over a field $\mathbb{K}$, and let
\[
f : V \rightarrow V
\]
be an endomorphism, that is, a linear map from $V$ to itself.

A nonzero vector $v \in V$ is called an eigenvector of $f$ if there exists a scalar
$\lambda \in \mathbb{K}$ such that
\[
f(v) = \lambda v.
\]
The scalar $\lambda$ is then called the eigenvalue associated with the eigenvector $v$.

This relation means that the map $f$ does not change the direction of the vector $v$.
It simply multiplies it by a scalar factor $\lambda$.

Now consider a basis
\[
\mathcal{B} = (e_1,\dots,e_N)
\]
of $V$. As recalled in Appendix~\ref{app:representation-lineaire-matrice},
the endomorphism $f$ can be represented by a matrix in this basis.

A natural question is whether we can choose a basis in which the matrix
of $f$ becomes particularly simple.

The most favorable case is when the basis $\mathcal{B}$ consists entirely of
eigenvectors of $f$. Suppose, therefore, that for every $i \in \{1,\dots,N\}$,
\[
f(e_i) = \lambda_i e_i.
\]

Then, in the basis $\mathcal{B}$, the coordinates of $f(e_i)$ are zero everywhere except at the
$i$-th position:
\[
[f(e_i)]_{\mathcal{B}}
=
\begin{pmatrix}
0 \\
\vdots \\
0 \\
\lambda_i \\
0 \\
\vdots \\
0
\end{pmatrix}.
\]

The matrix representing $f$ in this basis therefore has these coordinate vectors as its columns.
It is diagonal:
\[
A_{\mathcal{B}}
=
\begin{pmatrix}
\lambda_1 & 0         & \cdots & 0 \\
0         & \lambda_2 & \cdots & 0 \\
\vdots    & \vdots    & \ddots & \vdots \\
0         & 0         & \cdots & \lambda_N
\end{pmatrix}.
\]

In other words, given a basis of eigenvectors, the action of $f$ becomes very simple:
each coordinate is merely multiplied by the corresponding eigenvalue.

Conversely, if the matrix of $f$ is diagonal in a given basis, then the vectors of that
basis are eigenvectors of $f$.

This observation explains why it is useful to find a basis of eigenvectors: in such a basis,
the linear map no longer mixes directions. It acts independently on each direction
of the basis.

In this report, this idea is applied to circular convolution. The matrix associated
with circular convolution has the Fourier vectors as eigenvectors. Thus, when
the signal is expressed in the Fourier basis, the convolution operator becomes diagonal:
\[
C_h = F_N \Lambda_h F_N^{-1},
\]
where $F_N$ contains the Fourier vectors and $\Lambda_h$ is the diagonal matrix of the
associated eigenvalues.
